\section{Architecture}
\label{sec:architecture}

All models in Sections~\ref{sec:training-evolution}--\ref{sec:embedding-analysis}
share the same \emph{CNN-Transformer} architecture (referred to as
\texttt{CNNTransformerV4} in the codebase), operating on a $40 \times 40$
binary grid with toroidal (periodic) boundary conditions. The forward pass
has three stages:

\paragraph{Stage 1 --- Local CNN feature extractor.}
Two $3\times 3$ convolutions with \texttt{circular} padding mode (so that
convolution respects the toroidal topology exactly, rather than treating
the grid border as a hard edge), each followed by GroupNorm and a GELU
nonlinearity:
\[
  \text{Conv2d}(1 \to 32) \to \text{GroupNorm} \to \text{GELU}
  \to \text{Conv2d}(32 \to d_{\text{model}}) \to \text{GroupNorm} \to \text{GELU}.
\]
The receptive field of this stage is $5\times 5$, matching (and slightly
exceeding) the Moore neighborhood relevant to the GoL update rule.

\paragraph{Stage 2 --- Patch tokenization.}
The $40\times 40$ feature map is partitioned into non-overlapping
$4\times 4$ patches (a $10\times 10$ arrangement, $100$ patches total),
following the patch-tokenization scheme of Vision Transformers
\cite{dosovitskiy2020vit}. Each
patch's $d_{\text{model}} \times 4 \times 4$ feature block is flattened and
linearly projected to a $d_{\text{model}}$-dimensional token, and a
\emph{learnable} absolute positional embedding
$\texttt{pos\_embed} \in \mathbb{R}^{100 \times d_{\text{model}}}$ is added:
\[
  \text{token}_i = \text{Linear}(\text{feat}_i) + \text{pos\_embed}_i,
  \qquad i = 1, \dots, 100.
\]
We refer to this as the \textbf{pre-transformer embedding} in
Section~\ref{sec:embedding-analysis}.

\paragraph{Stage 3 --- Transformer encoder.}
A pre-norm Transformer encoder \cite{vaswani2017attention} (GELU
activation, $\text{dim\_feedforward} = 4 d_{\text{model}}$) mixes
information across all 100 patch tokens via
self-attention. We refer to its output as the \textbf{post-transformer
embedding}. A final linear head maps each output token back to
$4\times 4 = 16$ logits, reconstructing next-state predictions for every
cell in the patch.

\paragraph{Model sizes used.} Two sizes appear throughout this report:

\begin{center}
\begin{tabular}{lccccc}
\toprule
Name & $d_{\text{model}}$ & heads & layers & dim\_feedforward & params \\
\midrule
V8 (baseline size)  & 64  & 4 & 4 & 256 & 291{,}888 \\
V10 (scaled size)   & 128 & 8 & 6 & 512 & 1{,}504{,}240 \\
\bottomrule
\end{tabular}
\end{center}

\paragraph{Data augmentation.} Training uses the dihedral group $D_4$
(90$^\circ$ rotations and horizontal/vertical reflections) applied jointly
to input and target frames. Translation augmentation is deliberately
\emph{not} used, since it directly conflicts with the learnable absolute
positional embedding in Stage 2 (a translated patch would need a different
position embedding than the one it is assigned).
